\begin{figure}[tbp]
\centering
\begin{tikzpicture}[
  node distance=4mm and 8mm,
  box/.style={draw=linkblue!70,rounded corners=2pt,align=center,
    text width=5.4cm,minimum height=10mm,inner sep=5pt,font=\small},
  closure/.style={box,dashed,text width=4.2cm},
  arr/.style={-{Stealth[length=2mm]},draw=linkblue!80,thick}]
\node[box] (orders) {Unbounded categoricity\\Order presentations and orbit reversal};
\node[box,below=of orders] (diagrams) {Single-model diagrams\\Uniform testing and compact placement};
\node[box,below=of diagrams] (calculus) {Working schemes\\Free extension and relative isolation\\Extensions controlled by a mark};
\node[box,below=of calculus] (point) {Minimal type\\Fixed pointed models};
\node[box,below=of point] (marked) {Finite marked systems\\Exact boundary lifting\\at fixed depth};
\node[box,below=of marked] (growth) {No proper extension with unchanged points\\Basis size equals model size on a tail};
\node[box,below=of growth] (systems) {Geometric systems and relative bases\\Exact lifting for each finite shape};
\node[box,below=of systems] (tail) {Induction through larger cardinalities\\Categoricity on a tail};
\foreach \a/\b in {orders/diagrams,diagrams/calculus,calculus/point,
 point/marked,growth/systems}
  \draw[arr] (\a) -- (\b);
\node[closure,right=of marked] (markedclosure)
  {Common marked level\\All finite marked depths};
\draw[arr,dashed] (marked.east) -- (markedclosure.west);
\draw[arr,dashed] (markedclosure.south) |- (growth.east);
\node[closure,right=of systems] (geometricclosure)
  {Common geometric level\\All finite shapes\\Above the no-growth threshold};
\draw[arr,dashed] (systems.east) -- (geometricclosure.west);
\draw[arr,dashed] (geometricclosure.south) |- (tail.east);
\end{tikzpicture}
\caption{The qualitative proof, ordered by its principal dependencies.
The two dashed boxes are separate applications of
Section~\ref{close:section-uniform}. The first makes all finite marked
depths available for the no-growth argument over the fixed point. The
second is chosen after that argument
and supplies the starting cardinal for the final transfer. These
classes of common levels need not contain a tail. The small-code
reduction then gives the uniform threshold of
Theorem~\ref{main:uniform}.}
\label{fig:proof-route}
\end{figure}
